\documentclass[a4paper, serif, 10pt]{moderncv}

%% moderncv
% color and style
\moderncvcolor{black}
\moderncvstyle{banking}

%% page layout/margins
\usepackage[top=2.5cm, bottom=2.5cm, left=1.5cm, right=1.5cm]{geometry}

\nopagenumbers{}

%% Babel config for Spanish language
% \usepackage[spanish]{babel} has some conflicting issues with moderncv
% so we have to use enable the following options to make it work
%
% ref:
% - https://tex.stackexchange.com/a/140161/36007
\usepackage[spanish,es-lcroman]{babel}

%% fontspec
\usepackage{fontspec}

\IfFontExistsTF{Linux Libertine}{
  \setmainfont[Ligatures={TeX, Common}, Numbers=OldStyle]{Linux Libertine}
}{}
\IfFontExistsTF{Linux Libertine O}{
  \setmainfont[Ligatures={TeX, Common}, Numbers=OldStyle]{Linux Libertine O}
}{}

%% CTeX
% CJK support, used to show CJK characters in the resume
%
% - heading=false: disable ctex heading
% - punct=kaiming: use kaiming punctuations styles for CJK
% - scheme=plain: use plain scheme, do not override \normalsize font size
% - space=auto: space settings for CJK characters
%
% CJK fonts are resolved per locale language, e.g. Noto Serif/Sans CJK JP
% for Japanese resumes, with fallback to the next region if not installed
%
% ref:
% - http://ctan.mirrorcatalogs.com/language/chinese/ctex/ctex.pdf
\usepackage[UTF8, heading=false, punct=kaiming, scheme=plain, space=auto]{ctex}

\IfFontExistsTF{Noto Serif CJK SC}{
  \setCJKmainfont{Noto Serif CJK SC}
}{}
\IfFontExistsTF{Noto Sans CJK SC}{
  \setCJKsansfont{Noto Sans CJK SC}
}{}

%% line spacing
\usepackage{setspace}
\setstretch{1.125}

%% URL styling - use normal text instead of monospace
\urlstyle{same}

%% Patch \httplink and \httpslink to handle full URLs
% The original moderncv macros always prepend http:// or https:// to URLs.
% This causes issues when the URL already has a protocol (e.g., https://example.com)
% resulting in malformed URLs like https://https://example.com.
% This patch checks if the URL already contains "://" and uses it directly if so.
% Note: We use \str_set:Nx (x-expansion) to expand macros like \@homepage first.
\ExplSyntaxOn
\str_new:N \g_yamlresume_url_str

\RenewDocumentCommand{\httpslink}{O{}m}{
  \str_set:Nx \g_yamlresume_url_str {#2}
  \str_if_in:NnTF \g_yamlresume_url_str {://}
    { \url{#2} }
    { \url{https://#2} }
}

\RenewDocumentCommand{\httplink}{O{}m}{
  \str_set:Nx \g_yamlresume_url_str {#2}
  \str_if_in:NnTF \g_yamlresume_url_str {://}
    { \url{#2} }
    { \url{http://#2} }
}
\ExplSyntaxOff

\name{Carlos Javier Pérez Gómez}{}
\title{Ingeniero Mecánico Senior especializado en diseño de producto y fabricación}
\phone[mobile]{+34 678 123 456}
\email{carlos.perez@example.com}
\homepage{https://www.carlosperez-ingenieria.com}

\address{Calle de Alcalá, 123, 4ºB, Madrid, Comunidad de Madrid, España, 28001}{}{}

\extrainfo{{\small \faLinkedin}\ \href{https://www.linkedin.com/in/carlosperezmecanico}{@carlosperezmecanico} {} {} {} • {} {} {} 
{\small \faGithub}\ \href{https://github.com/cperez-mec}{@cperez-mec}}

\begin{document}

\maketitle

\section{Información básica}

\cvline{}{\begin{itemize}
\item Ingeniero Mecánico con más de 8 años de experiencia en el diseño, desarrollo y optimización de productos y sistemas mecánicos.
\item Especializado en CAD/CAE (SolidWorks, CATIA, ANSYS) y en la gestión de proyectos de fabricación.
\item Sólida experiencia en el sector de automoción y energías renovables, liderando equipos multidisciplinares.
\item Enfoque en la mejora continua, reducción de costes y cumplimiento de normativas de calidad (ISO 9001, IATF 16949).
\end{itemize}}
\section{Educación}

\cventry{sept 2012–jul 2016}
        {Licenciatura, Ingeniería Mecánica, Puntuación: 8.2/10}
        {Universidad Politécnica de Madrid}
        {\url{https://www.upm.es}}
        {}
        {\begin{itemize}
\item Formación sólida en fundamentos de ingeniería mecánica, con especial énfasis en diseño y análisis.
\item Participación en proyectos de fin de grado sobre optimización de estructuras ligeras.
\item Prácticas en laboratorio de ensayos mecánicos y caracterización de materiales.
\end{itemize}
\textbf{Cursos}: Cálculo Diferencial e Integral, Mecánica de Sólidos, Termodinámica Aplicada, Mecánica de Fluidos, Diseño de Máquinas, Materiales de Ingeniería}

\cventry{sept 2016–jul 2018}
        {Maestría, Ingeniería Mecánica Avanzada, Puntuación: 9.1/10}
        {Universidad Carlos III de Madrid}
        {\url{https://www.uc3m.es}}
        {}
        {\begin{itemize}
\item Especialización en simulación numérica y diseño avanzado de componentes mecánicos.
\item Trabajo de fin de máster sobre análisis de fatiga en palas de aerogeneradores.
\item Colaboración con departamento de I+D de empresa del sector eólico.
\end{itemize}
\textbf{Cursos}: Diseño Mecánico Asistido por Ordenador, Mecánica de Materiales Compuestos, Dinámica de Sistemas Multicuerpo, Tribología, Optimización Topológica, Simulación CFD}
\section{Experiencia laboral}

\cventry{sept 2018–dic 2021}
        {Ingeniero Mecánico de Diseño}
        {Siemens Gamesa Renewable Energy}
        {\url{https://www.siemensgamesa.com}}
        {}
        {\begin{itemize}
\item Diseño y desarrollo de componentes mecánicos para aerogeneradores de gran potencia.
\item Análisis estructural y de fatiga mediante ANSYS, optimizando la vida útil de los componentes.
\item Coordinación con equipos de fabricación y proveedores para asegurar la viabilidad del diseño.
\item Redacción de especificaciones técnicas y documentación de ingeniería.
\end{itemize}
\textbf{Palabras clave}: Aerogeneradores, Diseño mecánico, CATIA V5, ANSYS, Fatiga, Normativa IEC}

\cventry{ene 2022 hasta la fecha}
        {Ingeniero de Producto - Componentes Mecánicos}
        {SEAT S.A.}
        {\url{https://www.seat.es}}
        {}
        {\begin{itemize}
\item Gestión del ciclo de vida de componentes mecánicos para vehículos eléctricos.
\item Aplicación de GD\&T y normas IATF 16949 en la definición de especificaciones.
\item Liderazgo de equipos multidisciplinares en la resolución de problemas de calidad.
\item Implementación de acciones de reducción de costes sin comprometer el rendimiento.
\end{itemize}
\textbf{Palabras clave}: Automoción, SolidWorks, GD\&T, IATF 16949, APQP, Reducción de costes}
\section{Idiomas}

\cvline{Español}{Competencia nativa o bilingüe \hfill \textbf{Palabras clave}: Nativo}
\cvline{Inglés}{Competencia profesional plena \hfill \textbf{Palabras clave}: TOEFL iBT 110, Uso profesional}
\cvline{Francés}{Competencia limitada de trabajo \hfill \textbf{Palabras clave}: DELF B1}
\section{Competencias}

\cvline{Diseño Mecánico CAD}{Experto \hfill \textbf{Palabras clave}: SolidWorks, CATIA V5, AutoCAD, Inventor, GD\&T}
\cvline{Simulación y Análisis CAE}{Avanzado \hfill \textbf{Palabras clave}: ANSYS, Abaqus, MATLAB, Simulink, CFD}
\cvline{Fabricación y Procesos}{Avanzado \hfill \textbf{Palabras clave}: Lean Manufacturing, Six Sigma, CNC, Impresión 3D, Soldadura}
\cvline{Gestión de Proyectos}{Intermedio \hfill \textbf{Palabras clave}: PMP, Agile, Scrum, Microsoft Project, SAP}
\cvline{Programación y Herramientas}{Intermedio \hfill \textbf{Palabras clave}: Python, VBA, MATLAB, PLC, Arduino}
\section{Premios}

\cventry{nov 2019}
        {Colegio Oficial de Ingenieros Industriales de Madrid}
        {Premio a la Innovación en Diseño Mecánico}
        {}
        {}
        {Reconocimiento por el desarrollo de un sistema de transmisión más eficiente para aerogeneradores.}
\section{Certificados}

\cventry{mar 2018}
        {Dassault Systèmes}
        {Certified SolidWorks Professional (CSWP)}
        {\url{https://www.solidworks.com/certification}}
        {}
        {}

\cventry{jun 2020}
        {Project Management Institute}
        {Project Management Professional (PMP)}
        {\url{https://www.pmi.org/certifications/project-management-pmp}}
        {}
        {}

\cventry{sept 2019}
        {ASQ}
        {Lean Six Sigma Green Belt}
        {\url{https://asq.org/cert/lean-six-sigma-green-belt}}
        {}
        {}
\section{Publicaciones}

\cventry{dic 2021}
        {Análisis de fatiga en palas de aerogeneradores de gran potencia}
        {Revista Iberoamericana de Ingeniería Mecánica}
        {\url{https://www.revista-ingenieria-mecanica.org/articulo/12345}}
        {}
        {Estudio sobre la vida a fatiga de materiales compuestos en palas de aerogeneradores mediante simulación numérica y validación experimental.}
\section{Proyectos}

\cventry{ene 2020–dic 2020}
        {Proyecto de I+D para la optimización topológica de un chasis tubular de aluminio.}
        {Diseño de chasis ligero para vehículo eléctrico}
        {\url{https://www.proyectochasis.com}}
        {}
        {\begin{itemize}
\item Rediseño del chasis reduciendo un 18\% el peso manteniendo la rigidez torsional.
\item Validación mediante análisis FEM (ANSYS) y pruebas físicas en pista.
\item Coordinación con proveedores para la fabricación de prototipos.
\end{itemize}
\textbf{Palabras clave}: Optimización topológica, Aluminio, FEM, Vehículo eléctrico}
\section{Intereses}

\cvline{Automoción}{Vehículos eléctricos, Clásicos, Motos}
\cvline{Energías Renovables}{Eólica, Solar, Almacenamiento}
\cvline{Deporte}{Ciclismo, Senderismo, Natación}
\section{Voluntariado}

\cventry{sept 2019–jun 2022}
        {Mentor de Jóvenes Ingenieros}
        {Fundación Ingeniería Sin Fronteras}
        {\url{https://www.isf.es}}
        {}
        {\begin{itemize}
\item Mentoría a estudiantes de ingeniería en el desarrollo de proyectos de cooperación.
\item Organización de talleres sobre diseño sostenible y energías renovables.
\item Colaboración en la evaluación de prototipos para comunidades rurales.
\end{itemize}}
    
\end{document}